

\section{matrix element of \texorpdfstring{$e^{-iqx}$}{e\textsuperscript{-iqx}} in the harmonic oscillator basis}

the matrix element $\langle n | e^{-iqx} | m \rangle$ for the one-dimensional quantum harmonic oscillator eigenfunctions $|n\rangle$ and $|m\rangle$ is given in closed form by:

\begin{align}
  \langle n | e^{-iqx} | m \rangle = \sqrt{\frac{n_<!}{n_>!}} \left(-\frac{i q }{\sqrt{2}\alpha }\right)^{n_> - n_<} e^{-\frac{q^2 }{4\alpha ^2}} \, L_{n_<}^{(n_> - n_<)}\!\left(\frac{q^2 }{2\alpha ^2}\right)\label{eq:osc-me-expiqx}
\end{align}

\noindent where:
\begin{itemize}
  \item $\alpha  = \sqrt{\frac{m\omega}{\hbar}}$ is the oscillator parameter $[\alpha]=L^{-1}$,
  \item $n_< = \min(n, m)$ and $n_> = \max(n, m)$,
  \item $L_k^{(\rho)}(z)$ is the associated (generalized) Laguerre polynomial.
\end{itemize}

\subsection{derivation}

\subsubsection{1. ladder operator representation}
express the position operator $x$ in terms of the creation ($a^\dagger$) and annihilation ($a$) operators:
\begin{align}
  x = \frac{1}{\sqrt{2} \alpha } (a + a^\dagger)
\end{align}

defining the dimensionless parameter $\lambda = \dfrac{q }{\sqrt{2}\alpha }$, the exponent becomes $-i q x = -i \lambda (a + a^\dagger)$.

\subsubsection{2. operator disentangling (Baker--Campbell--Hausdorff)}
since $[a, a^\dagger] = 1$, the operator exponential factors as:
\begin{align}
  e^{-i \lambda (a + a^\dagger)} = e^{-i \lambda a^\dagger} e^{-i \lambda a} e^{-\frac{1}{2} \lambda^2}
\end{align}

\subsubsection{3. evaluation of matrix elements}
inserting this factorized form into the matrix element yields:
\begin{align}
  \langle n | e^{-iqx} | m \rangle = e^{-\frac{\lambda^2}{2}} \langle n | e^{-i \lambda a^\dagger} e^{-i \lambda a} | m \rangle
\end{align}

expanding both exponentials in power series:
\begin{align}
  e^{-i \lambda a} |m\rangle          & = \sum_{k=0}^{m} \frac{(-i\lambda)^k}{k!} \sqrt{\frac{m!}{(m-k)!}} |m-k\rangle  \\
  \langle n| e^{-i \lambda a^\dagger} & = \sum_{l=0}^{n} \frac{(-i\lambda)^l}{l!} \sqrt{\frac{n!}{(n-l)!}} \langle n-l|
\end{align}

taking the inner product $\langle n-l | m-k \rangle = \delta_{n-l, m-k}$ sets $l = k + n - m$. Assuming $n \ge m$ (so $n_> = n$ and $n_< = m$):
\begin{align}
  \langle n | e^{-iqx} | m \rangle = e^{-\frac{\lambda^2}{2}} \sum_{k=0}^{m} \frac{(-i\lambda)^k}{k!} \frac{(-i\lambda)^{k + n - m}}{(k + n - m)!} \frac{\sqrt{n!\, m!}}{(m - k)!}
\end{align}

factoring out $(-i\lambda)^{n-m}$:
\begin{align}
  \langle n | e^{-iqx} | m \rangle = e^{-\frac{\lambda^2}{2}} (-i\lambda)^{n-m} \sqrt{\frac{m!}{n!}} \sum_{k=0}^{m} \frac{(-1)^k \lambda^{2k}}{k! (m-k)!} \frac{n!}{(k + n - m)!}
\end{align}

\subsubsection{4. connection to Laguerre polynomials}
using the explicit series definition of the associated Laguerre polynomial,
\begin{align}
  L_m^{(\rho)}(z) = \sum_{k=0}^{m} (-1)^k \binom{m+\rho}{m-k} \frac{z^k}{k!} = \sum_{k=0}^{m} \frac{(-1)^k z^k}{k! (m-k)!} \frac{(m+\rho)!}{(k+\rho)!}
\end{align}

with $\rho = n - m$ and $z = \lambda^2 = \dfrac{q^2 }{2\alpha^2}$, the sum simplifies to:
\begin{align}
  \langle n | e^{-iqx} | m \rangle = \sqrt{\frac{m!}{n!}} \left(-\frac{i q }{\sqrt{2}\alpha }\right)^{n-m} e^{-\frac{q^2 }{4\alpha ^2}} L_m^{(n-m)}\!\left(\frac{q^2 }{2\alpha ^2}\right)
\end{align}

extending to arbitrary $n$ and $m$ by defining $n_<$ and $n_>$ completes the proof.
\newpage

\section{Calculation of the matrix element ${\cal V}$ in the displaced harmonic oscillator basis}\label{sec:osc-V-matrix-element}

we need to calculate the matrix element
\begin{align}
  {\cal V}(n_1,n_2;P_1,P_2;a)= \langle n_1(a) P_1|H_I|n_2(a) P_2\rangle
\end{align}
for the interaction Hamiltonian
\begin{align}
  H_I(y) = V_0 e^{-\frac{y^2}{2\sigma_V^2}}
\end{align}
which will be used for $y = X-x$.

by direct calculation we have
\begin{align}
  \langle n_1(a), P_1|H_I|n_2(a),P_2\rangle & = \frac{V_0}{2\pi\hbar}\int \limits_{-\infty}^{\infty} dx\int \limits_{-\infty}^{\infty} dX e^{-\frac{i}{\hbar}X(P_1-P_2)}u_{n_1}(x-a)u_{n_2}(x-a)e^{-\frac{(X-x)^2}{2\sigma_V^2}}\nonumber                         \\
  & =\frac{V_0}{2\pi\hbar}\int \limits_{-\infty}^{\infty} dx\int \limits_{-\infty}^{\infty} dX e^{-\frac{i}{\hbar}(X-x)(P_1-P_2)}e^{-\frac{i}{\hbar}x(P_1-P_2)}u_{n_1}(x-a)u_{n_2}(x-a)e^{-\frac{(X-x)^2}{2\sigma_V^2}}
\end{align}
after the change of variable $X-x=y$ the previous matrix element becomes
\begin{align}
  \langle n_1(a), P_1|H_I|n_2(a),P_2\rangle & = \frac{V_0}{2\pi\hbar}\int \limits_{-\infty}^{\infty} dy e^{-\frac{i}{\hbar}y(P_1-P_2)}e^{-\frac{y^2}{2\sigma_V^2}}\int \limits_{-\infty}^{\infty} dx e^{-\frac{i}{\hbar}x(P_1-P_2)}u_{n_1}(x-a)u_{n_2}(x-a)
\end{align}
The integral over $y$ can be calculated using (\ref{eq:app-gaussian-fourier-integral}), with $\lambda = \frac1{2\sigma_V^2}$ and $k = \frac{P_1-P_2}{\hbar}$
\begin{align}
  \int \limits_{-\infty}^{\infty} dy e^{-\frac{i}{\hbar}y(P_1-P_2)}e^{-\frac{y^2}{2\sigma_V^2}} = \sqrt{{2\pi\sigma_V^2}} e^{-\frac{(P_1-P_2)^2\sigma_V^2}{2\hbar^2}}
\end{align}
The integral over $x$ is calculated after a further change of variable $x-a=z$
\begin{align}
  \int \limits_{-\infty}^{\infty} dx e^{-\frac{i}{\hbar}x(P_1-P_2)}u_{n_1}(x-a)u_{n_2}(x-a) = e^{-\frac{i}{\hbar}a(P_1-P_2)}\int \limits_{-\infty}^{\infty} dz e^{-\frac{i}{\hbar}z(P_1-P_2)}u_{n_1}(z)u_{n_2}(z)
\end{align}
The remaining integral can be computed using (\ref{eq:osc-me-expiqx}), with $q = \frac{P_1-P_2}{\hbar}$ and $\alpha = \sqrt{\frac{m\omega}{\hbar}}$
\begin{align}
  \int \limits_{-\infty}^{\infty} dz e^{-\frac{i}{\hbar}z(P_1-P_2)}u_{n_1}(z)u_{n_2}(z) = \sqrt{\frac{n_<!}{n_>!}} \left(-\frac{i (P_1-P_2)  }{\sqrt{2}\hbar\alpha}\right)^{n_> - n_<} e^{-\frac{(P_1-P_2)^2 }{4\hbar^2\alpha^2}} \, L_{n_<}^{(n_> - n_<)}\!\left(\frac{(P_1-P_2)^2 }{2\hbar^2\alpha^2}\right)
\end{align}
and the final result is
\begin{importantbox}
  \begin{align}
    {\cal V}(n_1,n_2;P_1,P_2;a)&= \langle n_1(a) P_1|H_I|n_2(a) P_2\rangle  \nonumber\\
    &= \frac{V_0}{2\pi\hbar}\sqrt{{2\pi\sigma_V^2}} e^{-\frac{i}{\hbar}a(P_1-P_2)}e^{-\frac{(P_1-P_2)^2\sigma_V^2}{2\hbar^2}}\sqrt{\frac{n_<!}{n_>!}} \left(-\frac{i (P_1-P_2)  }{\sqrt{2}\hbar\alpha}\right)^{n_> - n_<} e^{-\frac{(P_1-P_2)^2 }{4\hbar^2\alpha^2}} \, L_{n_<}^{(n_> - n_<)}\!\left(\frac{(P_1-P_2)^2 }{2\hbar^2\alpha^2}\right)\label{eq:osc-V-n1n2}
  \end{align}
  where $n_<$ and $n_>$ are smaller/larger between $n_1$ and $n_2$
  For $P_1=P_2$, the expression containing the power of $P_1-P_2$
  should be evaluated through its limiting value.  Using the
  orthonormality of the harmonic-oscillator eigenfunctions gives
  \begin{align}
    {\cal V}(n_1,n_2;P,P;a)
    = \frac{V_0}{2\pi\hbar}\sqrt{2\pi\sigma_V^2}\,
    \delta_{n_1n_2}.\label{eq:osc-V-n1n2-diag}
  \end{align}
\end{importantbox}
